\section{Circuit Complexity --- PPoly}
%%% generated by the blueprint pipeline from TCSlib.Complexity.CircuitComplexity.PPoly
\subsection*{Circuit Complexity --- PPoly: Introduction}

Declarations of the Lean module \texttt{TCSlib.Complexity.CircuitComplexity.PPoly}.

\subsection*{Circuit Complexity --- PPoly: Results}

\begin{definition}[Languages recognized within a size bound]
\lean{Language.InSIZE}
\label{Language.InSIZE}
\leanok
\uses{BoolCircuit.DAGCircuit.size, BoolCircuit.DAGCircuitFamily, BoolCircuit.DAGCircuitFamily.HasFaninTwo, BoolCircuit.DAGCircuitFamily.language}
A language $L$ over Boolean strings belongs to $	extsf{SIZE}(T)$ if there is a family of
Boolean circuits, one for each input length, whose gates have fan-in at most two, whose
circuit at every length $n$ has size at most $T(n)$, and whose accepted language is
exactly $L$.
\end{definition}

\begin{definition}[Languages decidable in polynomial size]
\lean{Language.InPPoly}
\statementsource{Vad12}{pdf-bf2ce5422088-p058-b003}
\statementsource{arora-barak-ch07-randomized-computation}{
  pdf-26d95080dbde-p012-b003,
  pdf-26d95080dbde-p013-b001
}
\label{Language.InPPoly}
\leanok
\uses{Language.InSIZE}
A language $L$ over the Boolean alphabet is decidable by circuits of polynomial size if
there are natural numbers $a$ and $k$ such that $L$ has circuit complexity at most the
function $n\mapsto a(n+1)^k$.
\end{definition}

\begin{theorem}[Monotonicity of circuit size bounds]
\lean{Language.InSIZE.mono}
\label{Language.InSIZE.mono}
\leanok
\uses{Language.InSIZE}
\difficulty{4}
For functions $T,T' : \mathbb{N} \to \mathbb{N}$ and a language $L$ over Boolean
strings, if $L$ is recognized by a family of fan-in-two Boolean circuits whose size on
inputs of length $n$ is at most $T(n)$ for every $n$, and $T(n) \leq T'(n)$ for every
$n$, then $L$ is recognized by a family of fan-in-two Boolean circuits whose size on
inputs of length $n$ is at most $T'(n)$ for every $n$.
\end{theorem}

\begin{theorem}[A size bound implies polynomial size]
\lean{Language.InSIZE.inPPoly}
\label{Language.InSIZE.inPPoly}
\leanok
\uses{Language.InPPoly, Language.InSIZE, Language.InSIZE.mono}
\difficulty{1}
Let $T : \mathbb{N} \to \mathbb{N}$ and let $L$ be a language of Boolean strings. If $L$
is recognized by a well-formed circuit family with fan-in at most two and size at most
$T(n)$ on inputs of length $n$, and there are natural numbers $a,k$ such that $T(n) \le
a(n+1)^k$ for every $n$, then $L$ has polynomial circuit size with fan-in at most two.
\end{theorem}

\begin{theorem}[Polynomial-size circuit characterization]
\lean{Language.inPPoly_iff}
\label{Language.inPPoly_iff}
\leanok
\uses{BoolCircuit.DAGCircuitFamily, BoolCircuit.DAGCircuitFamily.HasFaninTwo, BoolCircuit.DAGCircuitFamily.IsPolySize, BoolCircuit.DAGCircuitFamily.language, Language.InPPoly}
\difficulty{4}
For every language $L$ over the Boolean alphabet, $L$ can be recognized by a family of
Boolean circuits of fan-in at most two whose size is bounded by a polynomial in the
input length if and only if there exists such a family whose recognized language is
exactly $L$.
\end{theorem}

\begin{theorem}[Positive Circuit Size]
\lean{BoolCircuit.DAGCircuit.one_le_size}
\label{BoolCircuit.DAGCircuit.one_le_size}
\leanok
\uses{BoolCircuit.DAGCircuit, BoolCircuit.DAGCircuit.size}
\difficulty{2}
Let $C$ be a Boolean circuit with $n$ inputs and gates in topological order, where each
gate at position $i$ reads only vertices indexed below $n+i$ and the output index is
below $n$ plus the number of gates. If the size of $C$ is the number of inputs plus the
number of gates, then its size is at least $1$.
\end{theorem}

\begin{theorem}[Evaluation of a One-Input Circuit]
\lean{BoolCircuit.DAGCircuit.eval_of_size_le_one}
\label{BoolCircuit.DAGCircuit.eval_of_size_le_one}
\leanok
\uses{BoolCircuit.DAGCircuit, BoolCircuit.DAGCircuit.eval, BoolCircuit.DAGCircuit.size, BoolCircuit.DAGCircuit.values_getD_input}
\difficulty{4}
Let $C$ be a directed acyclic Boolean circuit with one input and size at most $1$. For
every Boolean input assignment $x$, the output of $C$ is $x(0)$.
\end{theorem}

\begin{theorem}[No polynomial-size circuit bound]
\lean{Language.not_inSIZE_pow}
\label{Language.not_inSIZE_pow}
\leanok
\uses{BoolCircuit.DAGCircuit.one_le_size, BoolCircuit.DAGCircuit.size, Language.InSIZE}
\difficulty{5}
For every natural number $c$ and every language $L$ of Boolean lists, there is no family
of acyclic Boolean circuits with fan-in at most two that recognizes $L$ and whose
circuit on inputs of length $n$ has size at most $n^c$ for every $n$.
\end{theorem}

\begin{theorem}[Emptiness of polynomial-size circuit languages]
\lean{Language.setOf_inSIZE_pow_eq_empty}
\label{Language.setOf_inSIZE_pow_eq_empty}
\leanok
\uses{Language.InSIZE, Language.not_inSIZE_pow}
\difficulty{3}
The set of languages of Boolean lists recognized by some family of acyclic Boolean
circuits with fan-in at most two and size at most $n^c$ for every input length $n$, for
some natural number $c$, is empty.
\end{theorem}

\begin{theorem}[Eventual polynomial size characterization]
\lean{Language.inPPoly_iff_eventually}
\label{Language.inPPoly_iff_eventually}
\leanok
\uses{BoolCircuit.DAGCircuit.size, BoolCircuit.DAGCircuitFamily, BoolCircuit.DAGCircuitFamily.HasFaninTwo, BoolCircuit.DAGCircuitFamily.language, Language.InPPoly}
\difficulty{5}
A Boolean language $L$ has polynomial-size Boolean circuit families with fan-in at most
two if and only if there are a natural number $c$ and a family of directed acyclic
Boolean circuits with fan-in at most two that recognizes exactly $L$ and has size at
most $n^c$ for every input length $n\ge 2$.
\end{theorem}

\begin{definition}[Polynomial-time decidable languages]
\lean{BoolCircuit.PPoly}
\statementsource{Vad12}{pdf-bf2ce5422088-p058-b003}
\statementsource{arora-barak-ch07-randomized-computation}{
  pdf-26d95080dbde-p012-b003,
  pdf-26d95080dbde-p013-b001
}
\label{BoolCircuit.PPoly}
\leanok
\uses{Language.InPPoly}
Let a language over the Boolean alphabet be a set of finite Boolean strings. The class
$	ext{PPoly}$ consists of those languages that belong to the complexity class $	ext{P}$
and are decidable in polynomial time.
\end{definition}

\begin{theorem}[Characterization of polynomial circuit size]
\lean{BoolCircuit.mem_PPoly_iff}
\label{BoolCircuit.mem_PPoly_iff}
\leanok
\uses{BoolCircuit.PPoly, Language.InPPoly}
\difficulty{1}
For every language $L$ over the Boolean alphabet, $L$ belongs to the class of languages
recognized by polynomial-size circuits with fan-in at most two if and only if there are
natural numbers $a$ and $k$ and such a circuit family recognizing exactly $L$ whose size
on inputs of length $n$ is at most $a(n+1)^k$.
\end{theorem}
